\scanpage{75}
Posant
\[
\omega_V=a_0\wedge\mathbin{-}\wedge a_n
\]
on a
\[
\delta(f_V)
=\underbrace{\beta_0^n\beta_1^{n-1}\cdots\beta_{n-1}}_B
\rho\omega_V^{\otimes-2}
=B\rho\omega_V^{\otimes-2}
\]
($f$ lisse sss $B$ inv. i.e. les $\beta_i$ inv
\quad$0\leq i\leq\unsure{1|n-1}$)

Considérons alors
\[
\lambda=f_V(\widetilde s_0)
=\deleted{\unsure{-}}\frac{\sigma'\rho'}{\rho^2}
=\left\{\begin{array}{ll}
(n\text{ pair avec }\rho=2\rho')\ (2\text{ inv!})
&+\dfrac{\sigma'}{4\rho'}\\
n\text{ impair }\rho=\rho'
&+\dfrac{\sigma'}{\rho'}
\end{array}\right.
\]

La hyperquadrique affine circonscrite au polyèdre régulier est donc donnée par
\[
f_V-\lambda=0
\]
Pour qu'elle soit \deleted{\unsure{affine}}\inserted{projectivement}
lisse il f et s. que les $\beta_i$ et $\sigma'$ soient inv
(c'est donc le cas \uline{birégulier})

(D) \uline{NB} Supposons donnés $V$, $f_V$ et $\varphi_G^V$,
représentation
\[
\deleted{\text{\unsure{orth.}}}\qquad G\longrightarrow O(f_V)
\]
$1^\circ$ transformant les $\tau_i$ en réflexions \uline{orth.}
et satisfaisant de plus $2^\circ)$ $3^\circ)$ plus haut.
\deleted{\unsure{Supposons de plus}}\inserted{Ceci implique}
\[
\deleted{\unsure{4^\circ)}}
\left[
\deleted{f|(D_0=\bigcap_{1\leq i\leq n}H_i)}
\inserted{f|L_0}\neq0\quad\text{en ttes fibres.}
\right.
\]
\margin{(i.e. $\tau_0$ est une réflexion str. orth. de $V$)
\quad$\longleftrightarrow$}
\margin{\deleted{\unsure{C}}
\unsure{Cette condition entraîne celles des autres ?}}

Alors ça provient donc \unsure{fppf} d'une \unsure{sit.} \uline{\unsure{li}}
comme plus haut. Il suffit de choisir une section \inserted{$a_0$} de $D_0$
telle que $f_V(a_0)=1$

\mspage{74}

